\documentclass[12pt,a4paper]{ctexart}
\usepackage{amsmath,amssymb}
\usepackage{booktabs}
\usepackage{array}
\usepackage{tabularx}
\usepackage{geometry}
\usepackage{graphicx}
\usepackage{float}
\geometry{left=2.4cm,right=2.4cm,top=2.5cm,bottom=2.5cm}
\setlength{\parindent}{2em}
\renewcommand{\arraystretch}{1.15}

\newcommand{\norm}[1]{\left\lVert #1 \right\rVert}
\newcolumntype{Y}{>{\raggedright\arraybackslash}X}

\begin{document}
\thispagestyle{empty}

\begin{center}
    {\zihao{2}\bfseries 实验报告}

    \vspace{0.7em}
    {\zihao{4}数值分析实验一}

    \vspace{1em}
    姓名：蔡斌 \qquad 学号：1024001338 \qquad 班级：\mbox{数学jd2401}
\end{center}

\section{实验目的}

\begin{enumerate}
    \item 掌握列主元高斯消去法求解线性方程组的基本步骤。
    \item 掌握高斯--赛德尔迭代法的构造方式，结合迭代结果分析其收敛特性。
    \item 通过条件数和误差估计，分析线性方程组的解对右端项扰动的敏感性。
\end{enumerate}

\section{实验题目}

求解线性方程组
\[
A x = b,
\quad
A=
\begin{pmatrix}
10&7&8&7\\
7&5&6&5\\
8&6&10&9\\
7&5&9&10
\end{pmatrix},
\quad
b=
\begin{pmatrix}
32\\23\\33\\31
\end{pmatrix}.
\]

将右端项扰动为
\[
\hat{b}=
\begin{pmatrix}
32.01\\22.99\\33.01\\30.99
\end{pmatrix},
\]
再求解方程组 $A\hat{x}=\hat{b}$。比较实际相对误差与误差估计量在 $1$ 范数和 $\infty$ 范数下的结果，并据此评估该问题数值结果的可靠性。

\section{实验原理}

本次实验主要采用列主元高斯消去法和高斯--赛德尔迭代法求解线性方程组，并结合条件数分析误差。

\subsection{列主元高斯消去法原理}

高斯消去法的基本思想是通过初等行变换将原方程组化为上三角方程组，再由回代得到未知量。列主元高斯消去法是在每一步消去时，从当前列中选取绝对值最大的元素作为主元，必要时交换两行后再继续消元。

\subsection{高斯--赛德尔迭代法原理}

高斯--赛德尔迭代法是把线性方程组改写成迭代格式，然后从一个初始值出发，不断计算新的近似解，直到结果满足要求为止。本题中通过实际迭代结果可以看到，迭代值逐步稳定并接近方程组的解。

\subsection{误差分析原理}

设原方程组解为 $x$，扰动后方程组解为 $\hat{x}$，记
\[
\delta x=\hat{x}-x,\qquad \delta b=\hat{b}-b,
\]
则可以利用条件数估计解对右端项扰动的敏感程度。条件数越大，说明解对微小扰动越敏感。

\section{实验步骤}

这次实验按照以下步骤进行：
\begin{enumerate}
    \item 写出原方程组和扰动后的方程组，明确系数矩阵 $A$、右端项 $b$ 与 $\hat{b}$；
    \item 对原方程组使用列主元高斯消去法，记录主元选取、消元过程和回代结果；
    \item 对扰动后的方程组使用相同的消元顺序进行求解，并计算 $\delta x$ 与 $\delta b$；
    \item 将原方程组改写为高斯--赛德尔迭代格式，取初值 $x^{(0)}=(0,0,0,0)^T$ 进行迭代计算；
    \item 比较高斯--赛德尔迭代法与直接法所得结果；
    \item 在 $1$ 范数和 $\infty$ 范数下计算实际相对误差与误差估计上界，并据此分析结果可靠性。
\end{enumerate}

\section{实验结果}

本节分别给出列主元高斯消去法和高斯--赛德尔迭代法的计算结果。

\subsection{列主元高斯消去法}

列主元高斯消去法的具体流程如下：
\begin{enumerate}
    \item 对第 $k$ 列在当前主对角线以下元素中选取绝对值最大的元素作为主元；
    \item 若主元不在第 $k$ 行，则交换对应两行；
    \item 利用主元将其下方元素逐个消去；
    \item 当矩阵化为上三角形式后，再通过回代得到方程组的解。
\end{enumerate}

\subsubsection{求解 $Ax=b$}

原增广矩阵为
\[
\left[
\begin{array}{cccc|c}
10&7&8&7&32\\
7&5&6&5&23\\
8&6&10&9&33\\
7&5&9&10&31
\end{array}
\right].
\]

第 $1$ 列主元为 $10$，不需要换行。消元后得到
\[
\left[
\begin{array}{cccc|c}
10&7&8&7&32\\
0&\frac{1}{10}&\frac{2}{5}&\frac{1}{10}&\frac{3}{5}\\
0&\frac{2}{5}&\frac{18}{5}&\frac{17}{5}&\frac{37}{5}\\
0&\frac{1}{10}&\frac{17}{5}&\frac{51}{10}&\frac{43}{5}
\end{array}
\right].
\]

第 $2$ 列在第 $2\sim4$ 行中绝对值最大的是 $\frac{2}{5}$，交换第 $2$、第 $3$ 行，再继续消元，得
\[
\left[
\begin{array}{cccc|c}
10&7&8&7&32\\
0&\frac{2}{5}&\frac{18}{5}&\frac{17}{5}&\frac{37}{5}\\
0&0&-\frac{1}{2}&-\frac{3}{4}&-\frac{5}{4}\\
0&0&\frac{5}{2}&\frac{17}{4}&\frac{27}{4}
\end{array}
\right].
\]

第 $3$ 列主元取 $\frac{5}{2}$，交换第 $3$、第 $4$ 行后继续消元，得到上三角方程组
\[
\left[
\begin{array}{cccc|c}
10&7&8&7&32\\
0&\frac{2}{5}&\frac{18}{5}&\frac{17}{5}&\frac{37}{5}\\
0&0&\frac{5}{2}&\frac{17}{4}&\frac{27}{4}\\
0&0&0&\frac{1}{10}&\frac{1}{10}
\end{array}
\right].
\]

回代可得
\[
x_4=1,\qquad x_3=1,\qquad x_2=1,\qquad x_1=1.
\]
因此
\[
x=
\begin{pmatrix}
1\\1\\1\\1
\end{pmatrix}.
\]

\subsubsection{求解 $A\hat{x}=\hat{b}$}

由于系数矩阵 $A$ 没有变化，所以消元次序和消元系数与前面完全相同，只是右端项不同。经过同样的消元后，有
\[
\left[
\begin{array}{cccc|c}
10&7&8&7&32.01\\
0&\frac{2}{5}&\frac{18}{5}&\frac{17}{5}&7.402\\
0&0&\frac{5}{2}&\frac{17}{4}&6.7325\\
0&0&0&\frac{1}{10}&0.079
\end{array}
\right].
\]

回代得到
\[
\hat{x}=
\begin{pmatrix}
1.82\\-0.36\\1.35\\0.79
\end{pmatrix}.
\]

于是
\[
\delta x=
\hat{x}-x=
\begin{pmatrix}
0.82\\-1.36\\0.35\\-0.21
\end{pmatrix},
\qquad
\delta b=
\hat{b}-b=
\begin{pmatrix}
0.01\\-0.01\\0.01\\-0.01
\end{pmatrix}.
\]

\subsection{高斯--赛德尔迭代法}

把原方程组写成迭代格式：
\begin{align*}
x_1^{(k+1)}&=\frac{32-7x_2^{(k)}-8x_3^{(k)}-7x_4^{(k)}}{10},\\
x_2^{(k+1)}&=\frac{23-7x_1^{(k+1)}-6x_3^{(k)}-5x_4^{(k)}}{5},\\
x_3^{(k+1)}&=\frac{33-8x_1^{(k+1)}-6x_2^{(k+1)}-9x_4^{(k)}}{10},\\
x_4^{(k+1)}&=\frac{31-7x_1^{(k+1)}-5x_2^{(k+1)}-9x_3^{(k+1)}}{10}.
\end{align*}

取初值 $x^{(0)}=(0,0,0,0)^T$。前两次迭代结果为
\begin{align*}
x^{(1)}&=(3.2,\ 0.12,\ 0.668,\ 0.1988)^T,\\
x^{(2)}&=(2.44244,\ 0.180184,\ 1.0590176,\ 0.34708416)^T.
\end{align*}

继续迭代可以发现，该方法对本题是收敛的，但收敛速度较慢。

取停止准则
\[
\norm{x^{(k)}-x^{(k-1)}}_{\infty}<10^{-8},
\]
则对于原方程组，迭代在第 $4062$ 步得到
\[
x^{(4062)}\approx
\begin{pmatrix}
1.00000195\\
0.99999679\\
1.00000080\\
0.99999952
\end{pmatrix},
\qquad
\norm{Ax^{(4062)}-b}_{\infty}\approx 6.05\times 10^{-8}.
\]

对扰动后的方程组，迭代格式只需把右端项改成 $\hat{b}$。在同样的停止准则下，第 $3786$ 步得到
\[
\hat{x}^{(3786)}\approx
\begin{pmatrix}
1.81999805\\
-0.35999679\\
1.34999920\\
0.79000048
\end{pmatrix},
\qquad
\norm{A\hat{x}^{(3786)}-\hat{b}}_{\infty}\approx 6.04\times 10^{-8}.
\]

两组结果都与列主元高斯消去法一致。前 $10$ 次迭代数据见附录。

\section{结果分析}

题目给出的误差估计式为
\[
\frac{\norm{\delta x}}{\norm{x}}
\le
\operatorname{cond}(A)\,
\frac{\dfrac{\norm{\delta A}}{\norm{A}}+\dfrac{\norm{\delta b}}{\norm{b}}}{1-\operatorname{cond}(A)\dfrac{\norm{\delta A}}{\norm{A}}}.
\]

本题中系数矩阵没有扰动，即 $\delta A=0$，于是上式化为
\[
\frac{\norm{\delta x}}{\norm{x}}
\le
\operatorname{cond}(A)\frac{\norm{\delta b}}{\norm{b}}.
\]

先求矩阵的逆：
\[
A^{-1}=
\begin{pmatrix}
25&-41&10&-6\\
-41&68&-17&10\\
10&-17&5&-3\\
-6&10&-3&2
\end{pmatrix}.
\]

由此可得
\[
\norm{A}_1=\norm{A}_{\infty}=33,\qquad
\norm{A^{-1}}_1=\norm{A^{-1}}_{\infty}=136,
\]
所以
\[
\operatorname{cond}_1(A)=\operatorname{cond}_{\infty}(A)=33\times136=4488.
\]

\subsection{$1$ 范数下的比较}

\[
\norm{x}_1=4,\qquad
\norm{\delta x}_1=2.74,
\]
因此实际相对误差为
\[
\frac{\norm{\delta x}_1}{\norm{x}_1}=\frac{2.74}{4}=0.685.
\]

另一方面，
\[
\norm{b}_1=119,\qquad
\norm{\delta b}_1=0.04,
\]
所以
\[
\frac{\norm{\delta b}_1}{\norm{b}_1}=\frac{0.04}{119}\approx 3.3613\times 10^{-4}.
\]

误差估计上界为
\[
\operatorname{cond}_1(A)\frac{\norm{\delta b}_1}{\norm{b}_1}
=4488\times\frac{0.04}{119}
\approx 1.508571.
\]

于是
\[
0.685<1.508571.
\]

\subsection{$\infty$ 范数下的比较}

\[
\norm{x}_{\infty}=1,\qquad
\norm{\delta x}_{\infty}=1.36,
\]
因此实际相对误差为
\[
\frac{\norm{\delta x}_{\infty}}{\norm{x}_{\infty}}=1.36.
\]

又因为
\[
\norm{b}_{\infty}=33,\qquad
\norm{\delta b}_{\infty}=0.01,
\]
所以
\[
\frac{\norm{\delta b}_{\infty}}{\norm{b}_{\infty}}=\frac{0.01}{33}\approx 3.0303\times 10^{-4}.
\]

误差估计上界为
\[
\operatorname{cond}_{\infty}(A)\frac{\norm{\delta b}_{\infty}}{\norm{b}_{\infty}}
=4488\times\frac{0.01}{33}=1.36.
\]

因此
\[
\frac{\norm{\delta x}_{\infty}}{\norm{x}_{\infty}}=1.36,
\]
与理论上界正好相等。

\subsection{可靠性评估}

从结果看，右端项的扰动只有 $10^{-2}$ 量级，但解的相对变化已经达到 $68.5\%$ 和 $136\%$，这说明该方程组对右端项误差非常敏感。

再结合条件数
\[
\operatorname{cond}_1(A)=\operatorname{cond}_{\infty}(A)=4488,
\]
可以看出该方程组的条件数较大，因此解对右端项的微小变化比较敏感。对于原方程组，列主元高斯消去法可以求得精确解；但当右端项中出现较小误差时，解就可能发生较大的变化。

因此，本题结果的主要特点是：右端项虽只发生了很小的变化，解却发生了较大的变化。另一方面，$\infty$ 范数下实际相对误差恰好达到误差估计上界，说明该误差估计与实际结果比较接近。

\section{实验结论}

\begin{enumerate}
    \item 列主元高斯消去法求得原方程组的解为
    \[
    x=(1,1,1,1)^T.
    \]
    \item 对扰动后的右端项，解为
    \[
    \hat{x}=(1.82,\ -0.36,\ 1.35,\ 0.79)^T.
    \]
    \item 高斯--赛德尔迭代法与直接法结果一致，说明该迭代法对本题是可行的，但收敛速度较慢。
    \item 在 $1$ 范数下，实际相对误差为 $0.685$，估计上界为 $1.508571$；在 $\infty$ 范数下，实际相对误差与估计上界都等于 $1.36$。
    \item 该问题的条件数较大，说明方程组的解对右端项扰动十分敏感。即使右端项只有很小变化，解也可能产生较大变化，因此这个方程组的结果可靠性较差。
\end{enumerate}

\appendix
\pagestyle{plain}
\markboth{}{}

\section{$Ax=b$ 的前 $10$ 次高斯--赛德尔迭代}

\scriptsize
\setlength{\tabcolsep}{4pt}
\noindent
\resizebox{\textwidth}{!}{%
\begin{tabular}{c l c c}
\toprule
$k$ & $x^{(k)}$ & $\norm{x^{(k)}-x^{(k-1)}}_{\infty}$ & $\norm{Ax^{(k)}-b}_{\infty}$ \\
\midrule
1 & $(3.200000,\ 0.120000,\ 0.668000,\ 0.198800)^T$ & 3.200000 & 7.575600 \\
2 & $(2.442440,\ 0.180184,\ 1.059018,\ 0.347084)^T$ & 0.757560 & 4.587418 \\
3 & $(1.983698,\ 0.204917,\ 1.277715,\ 0.459009)^T$ & 0.458742 & 2.706187 \\
4 & $(1.713079,\ 0.209421,\ 1.390776,\ 0.544436)^T$ & 0.270619 & 1.533999 \\
5 & $(1.559680,\ 0.203082,\ 1.440415,\ 0.610310)^T$ & 0.153400 & 0.813859 \\
6 & $(1.478294,\ 0.191581,\ 1.453138,\ 0.661580)^T$ & 0.081386 & 0.461432 \\
7 & $(1.440277,\ 0.178267,\ 1.445396,\ 0.701816)^T$ & 0.040236 & 0.362124 \\
8 & $(1.427625,\ 0.165034,\ 1.427245,\ 0.733625)^T$ & 0.031809 & 0.286279 \\
9 & $(1.429143,\ 0.152881,\ 1.404695,\ 0.758934)^T$ & 0.025309 & 0.227784 \\
10 & $(1.437973,\ 0.142269,\ 1.381219,\ 0.779187)^T$ & 0.023476 & 0.182275 \\
\bottomrule
\end{tabular}%
}

\normalsize

\section{$A\hat{x}=\hat{b}$ 的前 $10$ 次高斯--赛德尔迭代}

\scriptsize
\setlength{\tabcolsep}{4pt}
\noindent
\resizebox{\textwidth}{!}{%
\begin{tabular}{c l c c}
\toprule
$k$ & $\hat{x}^{(k)}$ & $\norm{\hat{x}^{(k)}-\hat{x}^{(k-1)}}_{\infty}$ & $\norm{A\hat{x}^{(k)}-\hat{b}}_{\infty}$ \\
\midrule
1 & $(3.201000,\ 0.116600,\ 0.670240,\ 0.196784)^T$ & 3.201000 & 7.555608 \\
2 & $(2.445439,\ 0.173313,\ 1.063555,\ 0.343336)^T$ & 0.755561 & 4.569380 \\
3 & $(1.988501,\ 0.194496,\ 1.284499,\ 0.453752)^T$ & 0.456938 & 2.688739 \\
4 & $(1.719627,\ 0.195371,\ 1.399699,\ 0.537847)^T$ & 0.268874 & 1.516385 \\
5 & $(1.567989,\ 0.185331,\ 1.451349,\ 0.602529)^T$ & 0.151638 & 0.795693 \\
6 & $(1.488420,\ 0.170066,\ 1.465949,\ 0.652719)^T$ & 0.079569 & 0.451715 \\
7 & $(1.452291,\ 0.152934,\ 1.459959,\ 0.691966)^T$ & 0.039246 & 0.353218 \\
8 & $(1.441603,\ 0.135839,\ 1.443445,\ 0.722858)^T$ & 0.030892 & 0.278029 \\
9 & $(1.445156,\ 0.119790,\ 1.422429,\ 0.747310)^T$ & 0.024452 & 0.220065 \\
10 & $(1.456087,\ 0.105253,\ 1.400400,\ 0.766753)^T$ & 0.022030 & 0.174987 \\
\bottomrule
\end{tabular}%
}

\normalsize

\end{document}
